\section{Messages as evidence-bearing interfaces}\label{sec:C21:messages}
A delegation message is an interface contract. Its payload should separate instructions, observations, inferences, proposed actions, and unresolved questions. Flattening them into one natural-language paragraph forces the receiving model to reconstruct authority and provenance. This is especially dangerous when a worker quotes untrusted material that contains instructions. The manager must not treat quoted source content as a delegated command merely because it arrived through an internal agent channel.

An evidence-bearing result contains a claim identifier, source references, the relevant observation, the transformation applied, uncertainty, and applicability conditions. For a code change it includes the artifact digest, tested revision, test selection, and known failures. For a research claim it includes paper version, result location, population, baseline, and limitation. The representation can be structured JSON or another schema; the important property is that integration can validate required fields and inspect the source behind them.

Schema validity does not establish semantic validity. A worker can provide a well-formed citation that does not support the associated claim. Therefore validation must be layered. Structural validation checks field presence and types. Referential validation resolves identifiers and versions. Evidential validation tests whether the cited observation supports the inference. Task validation checks whether the work product addresses the delegated objective. Conflating these stages makes a successful parser appear to be a scientific verifier.

Compression at handoff introduces an information bottleneck. Suppose a worker observes $O$, emits message $Z$, and the manager chooses an action from $Z$. If two possible observations require disjoint correct actions but induce the same $Z$ and the same other manager information, no rule on that information can be correct on both. Prior knowledge or a follow-up observation can change the information set; the claim is not that all omitted text is unrecoverable. The data-processing principle motivates preserving sufficient evidence rather than only conclusions. In practice exact sufficient statistics are unknown, so interfaces should expose provenance links and permit targeted follow-up retrieval.

A useful failure case is a worker reporting that ``the benchmark improved'' while omitting that the improvement used additional attempts or a different evaluator. The manager integrates the result as a method gain, and the final synthesis becomes incorrect even though neither agent invented a number. A required comparison field would have exposed the mismatch. This example illustrates why evidence loss at handoff is a distinct error class from hallucination and deserves separate measurement.

Message cost also affects scaling. For a round in which each of $n$ workers emits an $L$-bit message to every other worker, payload volume is $n(n-1)L$ before protocol overhead. A star that sends those reports once to a manager and broadcasts an $L_b$-bit summary uses $nL+nL_b$ bits, but holding $L_b$ fixed can discard task-relevant distinctions. The bound is an accounting model, not a measured speedup or a guarantee that the summary is sufficient. Hierarchical summaries reduce volume but introduce another compression layer and a potential single point of epistemic failure. Topic-based routing and shared artifact stores can reduce duplication, provided access control and version identity remain explicit. The appropriate topology depends on the dependency graph; a universal full-mesh conversation is rarely a neutral baseline.

\subsection{Learned latent channels and representation compatibility}\label{sec:C21:latent}
Text is only one communication substrate. LatComm's October 2026 preprint encodes a context once, without a question, into fixed-size hidden-state slots; a different frozen receiver reads them through learned cross-attention and prompt-dependent gates. Its two stages train the interface, then freeze the compressor while adapting reader-interface and gate parameters for question answering; both LLM backbones remain frozen~\cite{lin2026latcomm}. Interlat instead studies direct last-layer-state communication and learned compression, including heterogeneous sender--receiver configurations~\cite{du2026interlat}. These mechanisms require access below a typical text-only hosted API. They must not be inferred from ordinary agent handoff support.

To compare such systems, let $x\in\mathcal X$ be sender evidence, $q\in\mathcal Q$ a receiver request, and $y\in\mathcal Y$ the target. A general channel has an encoder $e_\phi:\mathcal X\rightarrow\mathbb R^{p\times s\times d}$, a transport or quantization map $Q_b$, and a receiver distribution $p_\psi(y\mid q,z)$ with $z=Q_b(e_\phi(x))$. Here $p,s,d\in\mathbb N_{>0}$ are channel-group count, slots per group and slot width (with $p=1$ for a single group), and $b\in\mathbb N_{>0}$ is the stored bits per scalar. Under a fixed task distribution and nonnegative predictive-distribution loss $\ell$ (such as log loss), one possible joint design objective is
\begin{equation}
 \min_{\phi,\psi}\;
 \mathbb E\!\left[\ell\{p_\psi(\cdot\mid q,z),y\}
 +\lambda_B B(z)+\lambda_T T(x,q,z)\right].
\end{equation}
The weights $\lambda_B,\lambda_T\geq0$ convert transmitted bits $B$ and end-to-end time $T$ to the chosen objective units. This is a comparison formulation, not the objective attributed to either paper. Frozen-backbone methods restrict the trainable parameters to adapters, compressors, or readers. A reconstruction objective and a task objective impose different information preferences; successful reconstruction is not sufficient for correct action, while task-only optimization can discard distinctions needed by future requests outside the training distribution.

The channel's nominal payload is $psdb$ bits before headers and error-control overhead, but token count is an inadequate measure of its cost. The sender may process every evidence token, the receiver attends to slots at selected layers, and each sequential agent can retain a growing cache. Count transmitted tensors, storage copies, dtype conversions, encoder FLOPs, additional receiver attention, and device transfers. A comparison that charges generated text tokens to a text team while excluding latent tensors and extra hidden-state computation gives the methods different accounting boundaries. Conversely, text may be repeatedly tokenized and prefetched when a packet is reusable; reuse count and cache policy must therefore be matched.

Compatibility is semantic as well as dimensional. A width-matching projection does not establish that a new receiver assigns the same meaning to an old packet. Replacing the sender checkpoint, changing selected layers, quantizing the packet, or retraining the decoder changes the channel. Record an encoder--packet--reader compatibility identity and test migration with frozen packets as well as regenerated packets. If only regenerated packets work, the system has demonstrated a new channel, not backward-compatible memory. A context packet also needs source identity, validity interval, permissions, and deletion dependencies outside its latent values; a dense vector cannot be assumed to carry an enforceable authority label.

The inspected LatComm implementation makes this boundary identifiable: \texttt{LatentBridge.compress} constructs the packet, \path{LatCommModel.encode_context} obtains its encoder input, \texttt{LatentInjector} inserts receiver-side transformations, and \texttt{PacketCache} manages stored packets~\cite{lin2026latcomm}. These callable definitions were read at a dated source revision; they were not executed here. Their existence establishes an implementation path, not numerical parity with the paper or robust compatibility after modification.

The packet shape and cache key matter in the released code. \path{LatentBridge.compress} returns one slot array per paired layer, so a context packet is $[p,s,d]$, not merely $[s,d]$. \texttt{PacketCache} selects that packet by a hash of context text and loads its stored dtype; this lookup does not itself validate an encoder--reader compatibility identity. The generation path computes prompt-conditioned gates and reuses them during decoding~\cite{cumb2026latcommcode}. These observations bound the implementation claim: a wrapper must bind cache namespace and packet generation to the selected checkpoints and layer pairs. A matching text hash, shape or dtype alone does not establish semantic compatibility, and storage at 16 bits does not imply every consumer-side copy remains 16-bit. No mismatched-cache benchmark outcome is asserted.

A discriminating experiment compares raw text, equal-byte text compression, a single-vector representation, and a multi-slot channel under the same sender evidence, receiver, task stream, and total compute envelope. Include unseen questions about a fixed document and new document types; a query-independent packet should be tested on questions unavailable when it was encoded. Swap in mismatched packets as a negative control, remove the channel, and vary packet precision. Measure action correctness, contradiction handling, calibration, and effect-policy compliance alongside bandwidth and latency distributions. For model migration, cross old and new encoder--reader pairs and retain the off-diagonal failures. A large speed gain on repeated multiple-choice reading would remain useful, but would not by itself establish reliable sharing of open-ended plans, revocable authority, or concurrent agent state.

